\chapter{Kiến trúc và lựa chọn công nghệ}
\label{chap:architecture}

\section{Kiến trúc giải pháp}

\subsection{Tổng quan kiến trúc}

Hình~\ref{fig:arch} mô tả kiến trúc đầu cuối gồm sáu tầng. Tầng nguồn gồm bốn nguồn dữ liệu; tầng tiếp nhận dùng Kafka làm event bus duy nhất, với Debezium bắt thay đổi từ WAL của PostgreSQL; tầng xử lý dùng Spark cho cả luồng và lô cùng các dịch vụ tiêu thụ sự kiện; tầng lưu trữ gồm Delta Lake, PostgreSQL cho mart thời gian thực BR4--BR5, Redis và OpenSearch; tầng phục vụ gồm Trino và API FastAPI; trên cùng là ứng dụng web cho năm nhóm người dùng.

\begin{figure}[H]
\centering
\begin{tikzpicture}[node distance=5mm and 3mm, font=\footnotesize]
  % ---- tầng nguồn ----
  \node[box, text width=3.0cm] (cli) {\textbf{Client web / mobile}\\\scriptsize telemetry JSON (xAPI)};
  \node[box, text width=3.0cm, right=of cli] (pg) {\textbf{LMS PostgreSQL}\\\scriptsize OLTP, WAL};
  \node[box, text width=3.0cm, right=of pg] (obj) {\textbf{MinIO}\\\scriptsize học liệu theo phiên bản\\(content-sync đọc khi có CDC)};
  \node[box, text width=3.0cm, right=of obj] (sb) {\textbf{Sandbox}\\\scriptsize nhật ký chấm bài};
  % ---- tầng tiếp nhận ----
  \node[stage, below=4mm of pg, text width=3.0cm] (dbz) {Debezium (pgoutput)};
  \coordinate (kc) at ($(cli.west)!0.5!(sb.east)$);
  \node[hl, anchor=north, text width=12.9cm] (kafka) at (kc |- {$(dbz.south)+(0,-0.45)$}) {\textbf{Apache Kafka} --- event bus, 26 topic: telemetry, \texttt{cdc.lms.*}, \texttt{lms.*}, \texttt{grading.*}, \texttt{risk.*}, \texttt{content.*}, \texttt{notification.*}};
  % ---- tầng xử lý ----
  \node[stage, anchor=north west, text width=4.0cm, minimum height=17mm] (sss) at ($(kafka.south west)+(0,-0.55)$) {\textbf{Spark Structured}\\\textbf{Streaming}\\\scriptsize landing bronze/silver (1--5 phút)\\ telemetry-live, lớp phủ BR7/BR8\\ cửa sổ rủi ro BR5};
  \node[stage, right=3mm of sss, text width=4.0cm, minimum height=17mm] (batch) {\textbf{Spark batch + Airflow}\\\scriptsize gold, mart, features\\ bảo trì OPTIMIZE / VACUUM\\ reindex học liệu};
  \node[stage, right=3mm of batch, text width=4.0cm, minimum height=17mm] (svc) {\textbf{Dịch vụ sự kiện}\\\scriptsize notification (BR1), grading (BR3)\\ content-sync (BR8)\\ br45 CDC worker (BR4/BR5)};
  % ---- tầng lưu trữ ----
  \node[box, anchor=north west, text width=5.5cm, minimum height=14mm, draw=gold!80!black, fill=gold!8] (delta) at ($(sss.south west)+(0,-0.55)$) {\textbf{Delta Lake trên MinIO + Hive Metastore}\\\scriptsize bronze $\rightarrow$ silver $\rightarrow$ gold (Kimball) $\rightarrow$ mart / features};
  \node[box, right=3mm of delta, text width=2.0cm, minimum height=14mm] (pga) {\textbf{PostgreSQL}\\\scriptsize \texttt{analytics}\\ BR4/BR5};
  \node[box, right=3mm of pga, text width=1.7cm, minimum height=14mm] (redis) {\textbf{Redis}\\\scriptsize cache, lớp phủ, dedup};
  \node[box, right=3mm of redis, text width=1.9cm, minimum height=14mm] (os) {\textbf{OpenSearch}\\\scriptsize BM25 + kNN};
  % ---- tầng phục vụ ----
  \node[stage, below=5mm of delta, text width=5.5cm] (trino) {\textbf{Trino} --- SQL warehouse, phân quyền theo vai trò};
  \node[stage, anchor=north west, text width=7.0cm] (api) at ($(pga.south west)+(0,-0.5)$) {\textbf{FastAPI} --- API có xác thực, phạm vi theo khoá học / tổ chức};
  % ---- người dùng ----
  \node[hl, anchor=north, text width=12.9cm] (web) at (kc |- {$(api.south)+(0,-0.5)$}) {\textbf{Ứng dụng web (React)} --- Learner \textbullet{} Instructor \textbullet{} Advisor \textbullet{} Org admin \textbullet{} Data platform};
  % ---- mũi tên ----
  \draw[arr] (cli.south) -- (cli.south |- kafka.north);
  \draw[arr] (pg) -- (dbz);
  \draw[arr] (dbz) -- (dbz.south |- kafka.north);
  \draw[arr] (sb.south) -- (sb.south |- kafka.north);
  \draw[arr] (sss.north |- kafka.south) -- (sss.north);
  \draw[arr] (svc.north |- kafka.south) -- ($(svc.north)$);
  \draw[arr] (sss.south -| delta.north) -- (delta.north -| sss.south);
  \draw[arr, <->] (batch.south) -- (batch.south |- delta.north);
  \draw[arr] (svc.south) -- (svc.south |- redis.north);
  \draw[arr] (delta) -- (trino);
  \draw[arr] (trino.east) -- (trino.east -| api.west);
  \draw[arr] (pga.south) -- (pga.south |- api.north);
  \draw[arr] (redis.south) -- (redis.south |- api.north);
  \draw[arr] (os.south) -- (os.south |- api.north);
  \draw[arr] (api.south) -- (api.south |- web.north);
  % ---- nhãn tầng ----
  \node[grplabel, rotate=90, anchor=south] at ($(cli.west)+(-0.25,0)$) {NGUỒN};
  \node[grplabel, rotate=90, anchor=south] at ($(kafka.west)+(-0.25,0)$) {TIẾP NHẬN};
  \node[grplabel, rotate=90, anchor=south] at ($(sss.west)+(-0.25,0)$) {XỬ LÝ};
  \node[grplabel, rotate=90, anchor=south] at ($(delta.west)+(-0.25,0)$) {LƯU TRỮ};
  \node[grplabel, rotate=90, anchor=south] at ($(trino.west)+(-0.25,0)$) {PHỤC VỤ};
\end{tikzpicture}
\caption{Kiến trúc tổng thể của nền tảng Project}
\label{fig:arch}
\end{figure}

Ba nguyên tắc xuyên suốt kiến trúc:
\begin{itemize}
  \item \textbf{Kafka là ranh giới giữa các context.} Các bounded context chỉ giao tiếp qua topic Kafka và bảng lakehouse, không dùng chung cơ sở dữ liệu. Nhờ đó mỗi context có thể triển khai và thay đổi độc lập.
  \item \textbf{Spark là bên ghi duy nhất của lakehouse} (ADR 0011). Các job thời gian thực chỉ ghi ra Kafka và Redis. Điều này tránh xung đột commit khi nhiều engine ghi cùng bảng Delta và giữ logic idempotent ở một nơi.
  \item \textbf{Lakehouse là nguồn sự thật cho phân tích.} Redis và các mart phục vụ đều dựng lại được từ lakehouse hoặc từ dữ liệu thô (ADR 0006).
\end{itemize}

\subsection{Lựa chọn công nghệ}

Bảng~\ref{tab:tech} liệt kê công nghệ theo vai trò. Thành phần chính gồm Kafka~\cite{kreps2011kafka}, Debezium~\cite{debezium}, Spark~\cite{zaharia2016spark} với Structured Streaming~\cite{armbrust2018structured}, Delta Lake~\cite{armbrust2020delta}, Trino (hậu duệ của Presto~\cite{sethi2019presto}), Redis~\cite{redis} và OpenSearch~\cite{opensearch}. Nhóm giữ các lựa chọn công nghệ ban đầu nếu chạy được trên máy cá nhân; khi thay thế, chọn thành phần cùng vai trò và ghi lại bằng ADR.

{\small
\begin{longtable}{@{}L{3.2cm}L{3.4cm}L{7.9cm}@{}}
\caption{Công nghệ theo vai trò}\label{tab:tech}\\
\toprule
\textbf{Vai trò} & \textbf{Công nghệ} & \textbf{Lý do chọn} \\
\midrule
\endfirsthead
\toprule
\textbf{Vai trò} & \textbf{Công nghệ} & \textbf{Lý do chọn} \\
\midrule
\endhead
Tiếp nhận, event bus & Apache Kafka & Bộ đệm cho telemetry tốc độ cao, giữ thứ tự theo khoá, phát lại được; dùng chung làm event bus cho các sự kiện LMS (ADR 0003) \\
CDC & Debezium (pgoutput) & Đọc replication slot nên không tạo tải truy vấn lên production; giữ LSN để sắp thứ tự \\
Xử lý luồng & Spark Structured Streaming & Watermark theo event time, checkpoint, sink giao dịch Delta; dùng chung engine với batch (ADR 0009, 0014) \\
Xử lý lô & Apache Spark 3.5 & Hệ sinh thái Delta, xử lý song song cho model SQL và chia đoạn học liệu (ADR 0004) \\
Lakehouse & Delta Lake 3.3 trên MinIO & ACID MERGE, time travel, Change Data Feed; tương thích Databricks (ADR 0005) \\
Catalog & Hive Metastore 3.1.3 & Danh mục chung cho Spark và Trino; bản 3.1 tương thích Spark 3.5 (ADR 0010) \\
SQL warehouse & Trino & Truy vấn trực tiếp bảng Delta, ANSI SQL, phân quyền theo schema, không cần bản sao thứ hai (ADR 0010) \\
Phục vụ độ trễ thấp & Redis & Tra cứu theo khoá dưới mili-giây, TTL; dùng cho cache, lớp phủ thời gian thực, dedup thông báo (ADR 0006) \\
Mart thời gian thực BR4/BR5 & PostgreSQL 16 & Phục vụ giao dịch, cập nhật đúng ngay khi nguồn sửa (ADR 0009) \\
Tìm kiếm & OpenSearch 2.17 & BM25 kết hợp vector HNSW; xếp hạng lai bằng reciprocal rank fusion (ADR 0013) \\
Điều phối & Apache Airflow & DAG theo lịch và theo dataset; mọi task idempotent theo ngày chạy \\
API, giao diện & FastAPI, React 18 + TypeScript & Một API chung cho mọi BR; giao diện không dùng thư viện UI hay biểu đồ ngoài \\
\bottomrule
\end{longtable}}

\subsection{Kiến trúc phần mềm: DDD và kiến trúc lục giác}

Mã nguồn nằm trong một gói Python duy nhất (monorepo, ADR 0002), chia thành bảy bounded context theo Domain-Driven Design~\cite{evans2003}: \code{telemetry}, \code{notification}, \code{recommendation}, \code{grading}, \code{risk}, \code{analytics} và \code{content}. Mỗi context gồm bốn lớp theo kiến trúc lục giác~\cite{cockburn2005} (Bảng~\ref{tab:layers}); một test kiến trúc kiểm tra quy tắc phụ thuộc giữa các lớp.

\begin{table}[H]
\centering\small
\caption{Các lớp trong một bounded context}\label{tab:layers}
\begin{tabularx}{\textwidth}{@{}L{2.8cm}YL{4.0cm}@{}}
\toprule
\textbf{Lớp} & \textbf{Chứa gì} & \textbf{Được phụ thuộc vào} \\
\midrule
Lớp domain & Thực thể, value object, aggregate, sự kiện miền, \en{policy} (ví dụ \code{EngagementPolicy}, \code{MinCohortPolicy}) & Chỉ \code{shared_kernel.domain} \\
Lớp application & Use case (mỗi lớp một hàm \code{execute}), các cổng \code{Protocol} & domain \\
Lớp infrastructure & Adapter: Kafka, Redis, Delta, S3, Trino, OpenSearch, YAML & Mọi thứ \\
Lớp entrypoints & Điểm ghép nối cho job Spark, dịch vụ, API & Mọi thứ \\
\bottomrule
\end{tabularx}
\end{table}

Lợi ích cốt lõi là \textbf{một logic, nhiều engine}. Use case \code{NormalizeEvents} (kiểm tra hợp lệ, chuyển schema, khử định danh, gắn cờ trễ) chạy trong job streaming, job silver của lakehouse và dịch vụ xử lý sự kiện, cho cùng kết quả về bút danh và quyết định đưa vào DLQ. Các ngưỡng nghiệp vụ trong SQL mart lấy từ value object của domain và được test đơn vị giữ đồng bộ.

\subsection{Ngữ nghĩa xử lý luồng}

\begin{itemize}
  \item \textbf{Thời gian sự kiện ở mọi nơi} theo mô hình Dataflow~\cite{akidau2015dataflow}. Watermark cho phép sự kiện đến muộn tối đa 5 phút với telemetry và 10 giây với CDC. Riêng cửa sổ rủi ro BR5 dùng watermark 21 ngày, vì dữ liệu học tập có thể được đồng bộ muộn.
  \item \textbf{Dữ liệu trễ không bị mất.} Sự kiện trễ hơn watermark vẫn được lưu và gắn cờ. Các model gold tự mở rộng cửa sổ tính toán về ngày của sự kiện trễ nhất (cơ chế \code{late_data}, Mục~\ref{sec:late}).
  \item \textbf{Đúng một lần.} Landing dùng checkpoint của Structured Streaming kết hợp sink giao dịch Delta. Lớp silver dùng MERGE idempotent theo khoá nghiệp vụ. Lớp phủ Redis ghi giá trị tuyệt đối thay vì cộng dồn, nên phát lại không đếm hai lần.
  \item \textbf{Thứ tự CDC.} Thay đổi được sắp theo \code{source.lsn} cho từng khoá. Thay đổi có LSN cũ hơn không thể ghi đè bản ghi mới hơn.
  \item \textbf{Đồng hồ thiết bị lệch.} Sự kiện có thời gian sau thời điểm tiếp nhận quá 2 phút được ghim về thời điểm tiếp nhận (\code{ClockSkewPolicy}). Nếu không, một thiết bị có đồng hồ chạy nhanh sẽ đẩy watermark lên và làm rơi sự kiện đúng hạn của mọi người khác.
\end{itemize}

\subsection{Các quyết định kiến trúc}

Mỗi quyết định quan trọng được ghi thành một ADR~\cite{nygard2011adr} trong bộ ADR của dự án. ADR đã chấp nhận không sửa; quyết định mới thay thế bằng ADR mới. Bảng~\ref{tab:adr} tóm tắt mười sáu ADR.

{\small
\begin{longtable}{@{}L{0.9cm}L{4.6cm}L{9.0cm}@{}}
\caption{Tóm tắt các quyết định kiến trúc (ADR)}\label{tab:adr}\\
\toprule
\textbf{ADR} & \textbf{Quyết định} & \textbf{Lý do và hệ quả chính} \\
\midrule
\endfirsthead
\toprule
\textbf{ADR} & \textbf{Quyết định} & \textbf{Lý do và hệ quả chính} \\
\midrule
\endhead
0001 & Ghi quyết định bằng ADR & Tránh tranh luận lại; lưu lý do, không chỉ kết quả \\
0002 & Monorepo Python, DDD + hexagonal & Use case độc lập engine; một chi phí nhỏ cho các lớp gián tiếp \\
0003 & Kafka + Debezium cho tiếp nhận và CDC & Giữ thứ tự theo thực thể; phải giám sát replication slot \\
0004 & Flink cho luồng, Spark cho lô & Lựa chọn ban đầu của nhóm; được điều chỉnh bởi ADR 0009 và 0014 \\
0005 & Lakehouse Delta Lake medallion & ACID upsert, time travel, ranh giới PII; cần job bảo trì chống file nhỏ \\
0006 & Redis cho tầng phục vụ & Dưới mili-giây, TTL; Redis không phải nguồn sự thật \\
0007 & Khử định danh ở ranh giới silver & HMAC-SHA256 có phiên bản khoá; báo cáo tổ chức áp dụng $k = 5$ \\
0008 & Runtime ban đầu trong tiến trình & Trình diễn đầu cuối trên laptop bằng adapter trong bộ nhớ; BR1--BR3 sau đó chuyển sang runtime thật (ADR 0016) \\
0009 & PostgreSQL mart + Spark Structured Streaming cho BR4/BR5 & Một driver Spark thay vì hoàn thiện cả hai engine; độ trễ CDC đến cảnh báo khoảng 10 giây \\
0010 & Trino + Hive Metastore làm SQL warehouse & Đóng vai ``Databricks SQL + catalog'' trên chính các bảng Delta, không có bản sao thứ hai \\
0011 & Spark là bên ghi duy nhất của Delta & Connector Flink cho Delta đã lỗi thời; tránh xung đột commit \\
0012 & Gold là lược đồ hình sao Kimball, khoá tất định & \code{xxhash64} của khoá tự nhiên; dựng lại tái tạo đúng khoá; thành viên suy luận cho chiều đến muộn \\
0013 & Đồng bộ nội dung BR8 bằng CDC vào OpenSearch & Một đường đồng bộ cho mọi bên ghi; idempotent theo phiên bản đã chỉ mục \\
0014 & Spark Structured Streaming cho telemetry nhanh và lớp phủ BR7/BR8 & Các job Flink chưa tồn tại; độ trễ đo được khoảng 7 giây đủ cho lớp phủ \\
0015 & Phiên đặc trưng BR2 trên Spark Structured Streaming + Redis & Spark 3.5 không hỗ trợ \code{update} cho \code{session_window}; state phiên nằm trong Redis, \code{SET NX} chống đếm lại \\
0016 & Runtime thật cho BR1--BR3 & Kafka, Redis, PostgreSQL thật; kết quả đến UI qua SSE; projector BR3 đồng bộ \code{public.*} \\
\bottomrule
\end{longtable}}

\subsection{Điều chỉnh so với kiến trúc tham chiếu}
\label{sec:deviations}

Nhóm giữ nguyên vai trò của kiến trúc tham chiếu nhưng điều chỉnh bốn điểm:

\begin{enumerate}
  \item \textbf{Spark Structured Streaming thay Flink cho BR5 và lớp phủ BR7--BR8} (ADR 0009, 0014). Kiến trúc tham chiếu dùng Flink~\cite{carbone2015flink} cho xử lý luồng có trạng thái; khi hiện thực, job Flink mới ở mức khung còn Spark Structured Streaming đã chạy cho landing lakehouse. Một engine giảm hạ tầng vận hành và cho phép tái dùng code chuẩn hoá silver. Đổi lại, độ trễ 7--10 giây, cao hơn mức dưới một giây của Flink, nhưng phù hợp với cảnh báo cho cố vấn và lớp phủ trên mart theo ngày. Các dịch vụ Flink vẫn có trong Compose để chuyển đổi khi cần.
  \item \textbf{Trino + Hive Metastore cho vai trò ``Databricks''} (ADR 0010). Delta Lake chỉ là định dạng bảng; Hive Metastore và Trino đảm nhận vai trò catalog và SQL warehouse của Databricks, chạy được cục bộ. Trên cloud có thể thay bằng Databricks SQL và Unity Catalog mà không đổi mô hình dữ liệu.
  \item \textbf{Redis thay Cassandra} cho tầng phục vụ (ADR 0006). Nhóm cân nhắc Redis và Cassandra; Redis đủ cho khối lượng hiện tại, còn adapter sau cổng cho phép thay khi dữ liệu vượt bộ nhớ.
  \item \textbf{BR2 dùng feature store hai nhánh.} Yêu cầu BR2 mô tả pipeline đặc trưng hai nhánh qua feature store; nhóm hiện thực cả hai: batch ghi \code{features.learner_longterm}, streaming giữ state phiên trong Redis. Gợi ý dùng feature làm ngữ cảnh và xếp hạng bằng luật nghiệp vụ (Mục~\ref{sec:br2}).
\end{enumerate}
